\documentclass[../../main]{subfiles}
\begin{document}

This chapter provides an overview of the theoretical concepts required for the remainder of this thesis.

The first section presents Constraint Programming, including the role of CP solvers and the MiniZinc modeling language used in the experimental evaluation.

The second section introduces Large Language Models, describing their probabilistic foundations, inference and decoding mechanisms, and prompting techniques.

\section{Constraint Programming}
CP is a paradigm for the declarative modeling and effective solution of large, particularly combinatorial, problems, especially in planning and scheduling~\cite{CPholygrail}.

A constraint can be thought of intuitively as a formalization of dependencies in real-world systems and their mathematical abstractions. A constraint is a logical relation among a set of variables, each taking a value in a given domain. The constraint thus restricts the possible values that variables can take; it represents partial information about the variables of interest. Constraints can also be heterogeneous, so they can bind variables from different domains, for example a length (number) with a word (string). The important feature of constraints is their declarative manner, i.e., they specify what relationship must hold without specifying a computational procedure to enforce that relationship~\cite{CPholygrail}.

Constraints arise naturally in most areas of human endeavor. They are the natural medium of expression for formalizing regularities that underlie the computational and physical world (whether natural or designed) and their mathematical abstractions. 

CP is the study of computational systems based on constraints. The idea of CP is to solve problems by stating constraints (requirements) about the problem area and, consequently, finding a solution satisfying all the constraints.

\subsection{Solvers}

Throughout this thesis, the term solver denotes a system that computes assignments satisfying a set of constraints and, when required, optimizes an objective function.

A ``Constraint Satisfaction Problem'' (CSP) can be formalized as a tuple
\[
\langle X, D, C \rangle,
\]
where \(X\) is a finite set of decision variables, \(D(x)\) denotes the domain associated with each variable \(x \in X\), and \(C\) is a set of constraints defined over subsets of \(X\). A solution is an assignment
\[
s : X \rightarrow \bigcup_{x \in X} D(x)
\]
such that \(s(x) \in D(x)\) for all \(x \in X\) and all constraints in \(C\) are satisfied~\cite{SeaPearl}.

The objective in a CSP is to determine whether at least one feasible assignment exists, or to enumerate all feasible assignments.

A ``Constraint Optimization Problem'' (COP) extends the CSP by introducing an objective function. Formally, a COP can be defined as
\[
\langle X, D, C, f \rangle,
\]
where \(f : S \rightarrow \mathbb{R}\) is an objective function defined over the set \(S\) of feasible assignments. The goal is to compute an optimal solution
\[
s^* = \arg\min_{s \in S} f(s)
\quad \text{or} \quad
s^* = \arg\max_{s \in S} f(s),
\]
depending on whether the problem is a minimization or maximization task.

Different solver families address CSPs and COPs using distinct mathematical abstractions and algorithmic techniques.

\subsubsection{Constraint Programming (CP) solvers}

Constraint Programming solvers operate directly on variables with finite or structured domains and heterogeneous constraints. Their fundamental mechanisms are constraint propagation, which removes inconsistent values from variable domains, and systematic search, typically implemented via backtracking guided by variable- and value-ordering heuristics. CP solvers natively support high-level global constraints equipped with specialized filtering algorithms that exploit combinatorial structure~\cite{CPvsMILPsolvers}. A representative example is Gecode~\cite{Gecode}.

\subsubsection{Portfolio Solvers}
Solving combinatorial search problems is computationally challenging, and numerous techniques and constraint solvers have been developed for this purpose. It has become clear that different solvers are better when solving different problem instances, even within the same problem class. It has also been shown that a single, arbitrarily efficient solver can be significantly outperformed by using a portfolio of possibly on-average slower solvers~\cite{PortfolioSolvers}.

Algorithm portfolios~\cite{AlgPortfolio} can be seen as instances of the more general Algorithm Selection problem~\cite{AlgSelProb} where, as reported in~\cite{AlgSelCSP}, the algorithm selection is performed case-by-case for each problem to solve. Within the context of constraint solving, a portfolio approach makes it possible to combine a number $m > 1$ of different constituent solvers $s_1, \dots, s_m$ in order to create a globally better constraint solver, dubbed a portfolio solver. When a new problem instance $p$ is encountered, the portfolio solver tries to predict the best constituent solver(s) $s_{i_1}, \dots, s_{i_k}$ (with $1 \leq i_j \leq m$ for $j = 1, \dots, k$) for solving $p$ and then runs them on $p$. Properly selecting and scheduling the solvers is a crucial step for the performance of a portfolio solver, and it is usually performed by exploiting Machine Learning techniques based on features extracted from the problem $p$ to solve.

In this work, starting from the idea behind portfolio solvers, we will try to evaluate the performance of a pre-trained LLM to understand if this could be a viable solution to the algorithm selection problem~\cite{AlgSelProb}, avoiding the cost of training for portfolio solvers, following the new paradigm proposed in~\cite{fromPFStoAS}.

\subsection{MiniZinc}
MiniZinc is an expressive modeling language for constraint programming that supports a wide range of solvers and represents a practical compromise among several design alternatives.

It arose from the need for a standard modelling language for constraint programming problems, and it also simplifies solver benchmarking~\cite{MiniZinc}.

For these reasons, all problems employed for testing in the course of this work were written using MiniZinc.

\subsubsection{A MiniZinc Example}
A MiniZinc problem specification has two parts: (a) the model, which describes
the structure of a class of problems; and (b) the data, which specifies one particular problem within this class. The pairing of a model with a particular data set is a model instance (sometimes abbreviated to instance).

The model and data may be in separate files. Data files can only contain assignments to parameters declared in the model. A user specifies data files on the command line, rather than naming them in the model file, so that the model file is not tied to any particular data file~\cite{MiniZinc}.

Each MiniZinc model is a sequence of items, which may appear in any order.
Consider the MiniZinc model and example data for a restricted job shop scheduling problem in \Cref{fig:minizincModelEx} and \Cref{fig:minizincDataEx}.

Line 0 is a comment, introduced by the ``\%'' character.

Lines 1--5 are ``variable declaration items''. Line 1 declares \texttt{size} to be an integer parameter, i.e., a variable that is fixed in the model. Line 20 (in the data file) is an ``assignment item'' that defines the value of \texttt{size} for this instance. Variable declaration items can include assignments, as in line 3. Line 4 declares \texttt{s} to be a 2D array of ``decision variables''. Line 5 is an integer variable with a restricted range. Decision variables are distinguished by the \texttt{var} prefix.

Lines 7--8 show a user-defined ``predicate item'', \texttt{no\_overlap}, which constrains two tasks given by start time and duration so that they do not overlap in time.

Lines 10--17 show a ``constraint item''. It uses the built-in \texttt{forall} to loop over each job, and ensure that: (line 12) the tasks are in order; (line 13) they finish before end; and (lines 14--16) that no two tasks in the same column overlap in time. Multiple constraint items are allowed, they are implicitly conjoined.

Line 19 shows a \texttt{solve item}. Every model must include exactly one solve
item. Here we are interested in minimizing the end time. We can also maximize
a variable or just look for any solution (\texttt{solve satisfy}).

There is one kind of MiniZinc item not shown by this example: ``include items''.
They facilitate the creation of multi-file models and the use of library files~\cite{MiniZinc}.

\begin{figure}[ht]
    \centering
    \includegraphics[width=0.75\linewidth]{minizincModelEx.png}
    \caption{MiniZinc model (\texttt{jobshop.mzn}) for the job shop problem (the example was taken from~\cite{MiniZinc}).}
    \label{fig:minizincModelEx}
\end{figure}
\begin{figure}[ht]
    \centering
    \includegraphics[width=0.75\linewidth]{minizincDataEx.png}
    \caption{MiniZinc data (\texttt{jobshop2x2.dzn}) for the job shop problem (the example was taken from~\cite{MiniZinc}).}
    \label{fig:minizincDataEx}
\end{figure}


\subsubsection{MiniZinc Challenge}
Comparing constraint programming systems is fraught with difficulty. The reason is that there are so many components to a modern CP system, only some of which are implemented by some systems. To claim that one CP system is ``better'' than another is a bold claim, since there is almost certainly some problem for which the ``worse'' system allows a stronger model, or a better search, and performs better.

MiniZinc makes a good attempt to handle the most obvious obstacle: there are hundreds of potential global constraints, most of which are handled by few or no systems. 
A standard input language enables direct comparison of different solvers.
Hence, every year since 2008 the MiniZinc Challenge has been comparing different solvers that support MiniZinc~\cite{PhilMznChallenge}.

The aim of the challenge is to compare various constraint solving technology on the same problems sets. The focus is on finite domain propagation solvers. An auxiliary aim is to build up a library of interesting problem models, which can be used to compare solvers and solving technologies.

Challenge participants provide a FlatZinc or MiniZinc solver and global constraint definitions specialized for their solver. Each solver is run on 100 MiniZinc model instances. For FlatZinc solvers, the MiniZinc compiler runs on the MiniZinc model and instance using the provided global constraint definitions to create a FlatZinc file. The resultant FlatZinc file is then given as input to the provided FlatZinc solver. For MiniZinc solvers, the MiniZinc model and data are input to the provided solver. 

Points are awarded for solving problems, speed of solution, and goodness of solutions (for optimization problems)~\cite{PhilMznChallenge}, however, the MiniZinc challenge scoring system is not actually utilized in this specific research, since it is a comparative scoring system, while an absolute one would be best suited for this specific research.

\section{Large Language Models}
\label{sec:LLM}
LLMs are statistical language models that estimate probability distributions over sequences of tokens in natural language by learning statistical patterns over large text corpora.

LLMs are the culmination of decades of progress in natural language processing (NLP) and machine learning research, and their development is largely responsible for the explosion of artificial intelligence advancements across the late 2010s and 2020s. Popular LLMs have become household names, bringing generative AI to widespread public attention. LLMs are also used widely in enterprises, with organizations investing heavily across numerous business functions and use cases~\cite{LLMsIBM}.

LLMs are built on a type of neural network architecture called a transformer~\cite{TransformerArchitecture}, which employs self-attention mechanisms to model dependencies among tokens in a sequence. Formally, given an input sequence \(x_1, \dots, x_n\), the model estimates the conditional distribution
\[
P(x_{n+1} \mid x_1, \dots, x_n),
\]
and text generation proceeds autoregressively by repeatedly sampling from this distribution.

Training is typically divided into a large-scale pretraining phase followed by alignment stages. During pretraining, the model learns general linguistic and semantic representations through self-supervised objectives. Alignment techniques, such as supervised fine-tuning and reinforcement learning from feedback, adapt the model to follow instructions, respect constraints, and produce outputs that better match user expectations~\cite{LLMsIBM}.

LLMs are, at their core, autoregressive statistical models that generate text by iteratively predicting the most probable subsequent token in a sequence, based on patterns learned from their training data.

\subsection{Inference and Decoding}

At inference time, the model produces logits \(l_k\) over the vocabulary, which are transformed into probabilities using a softmax function~\cite{InfeerenceTrainingLLMs}. Decoding strategies determine how the next token is selected from this distribution.

Modern LLMs typically generate text in a left-to-right, token-by-token fashion. For each prefix, the model computes a probability distribution of the next token over a fixed vocabulary. A decoding method defines how the generated token sequence is derived from these probability estimations. Deterministic approaches select the most probable token at each step, while stochastic approaches introduce controlled randomness~\cite{LLMDecoding}.

Notable deterministic decoding strategies include: ``Greedy Search'', which is limited to selecting the token with highest probability at each time step, or ``Beam Search''~\cite{BeamSearch}, which maintains a beam of the $k$ most probable sequences at each time step, where the hyperparameter $k$ is referred to as the beam width.

Some notable stochastic methods are: Temperature Sampling, which samples tokens from the estimated next-token distribution, with the skewness of the distribution controlled via a temperature hyperparameter $\tau$; and Top-p Sampling, which considers only the minimal set of most probable tokens covering a specified fraction $p$ of the distribution.

A central practical constraint is the context window, which limits the total number of tokens that can be processed jointly as input and output. This constraint directly influences prompt design, the amount of contextual information that can be provided, and the feasibility of maintaining conversational state.

\subsection{Relevance to Solver Selection}
In this research, the aim is to test the potential of LLMs when used for solver selection. In this context, prompts are engineered to enable the LLM to evaluate the characteristics of a problem instance, to make a decision that is as functional as possible. 

When applied to solver selection, an LLM functions as a decision component that maps problem representations to solver choices. Its performance depends on both theoretical aspects, such as its ability to model complex input-output relationships, and practical factors, including prompt formulation, context management, and API-level parameters.

\end{document}